\subsection{The action category as coefficient colimits}
\label{sec:action-colimits}
We prove Theorem~\ref{thm:action-colimits}. All colimits in this subsection
are taken in abelian groups.
\begin{proof}
The set $K(l,k)$ contains $0_Q$, is downward closed, and is closed under
finite joins: if $ql,q'l\leq k$, then
$(q\vee q')l=ql\vee q'l\leq k$. Hence it is directed.
Define a category with objects $L$ and prospective arrow groups
\[
 H(l,k)=\varinjlim_{q\in K(l,k)}F(q).
\]
Every element has one coefficient representative $[q,x]$, because a
finite sum can be transported to the join of its indices and added
there. Two representatives agree exactly when their images agree at a
common later index; this follows directly from the directed-colimit
relations.

Define composition by
\[
 [p,y]\,[q,x]=[pq,\mu_{p,q}(y,x)].
\]
Its index is valid: $ql\leq h$ and $ph\leq k$ imply $pql\leq k$.
Naturality of $\mu$ makes the definition independent of transported
representatives, and bilinearity follows after moving finite sums to
common indices. Associativity and the global unit follow from the lax
monoidal axioms. The unit coefficient in $F(1_Q)$ has valid index at
$(l,l)$ since $1_Ql=l$. If its class is zero, the resulting object is
correctly a zero object. Thus $H$ is a preadditive category.

Send a scalar generator $a_{q,x,l}$ to $[q,x]$, and a transition
$t_{l,k}$ to the class of the global unit at support $1_Q$.
Every defining relation of $\mathcal A_L$ holds: transitions compose
by the unit law, coefficient transitions are the colimit relations,
scalar-transition naturality uses the same coefficient at the same
index, and scalar products use the displayed composition. Additivity
at a fixed scalar support is group addition; at different supports it
follows by transport to their join. This gives a functor
$\mathcal A_L\to H$.

Conversely assign
\[
 [q,x]\longmapsto t_{ql,k}a_{q,x,l}\quad\text{in }\mathcal A_L(l,k).
\]
For $q\leq q'$ in $K(l,k)$, the coefficient-transition relation gives
\[
 t_{q'l,k}a_{q',F_{q,q'}(x),l}
 =t_{q'l,k}t_{ql,q'l}a_{q,x,l}
 =t_{ql,k}a_{q,x,l}.
\]
Thus the map descends to the colimit. Scalar-transition naturality and
scalar composition show that it preserves composition, with order
$pq$ and coefficient order $y,x$, without commuting them. The two
functors are inverse on all coefficient representatives and on the
scalar and transition generators. This proves the exact hom-group
identification, rather than only a surjection.

The poset $D_l$ contains zero and is directed by pointwise finite
joins. Such joins preserve $Q$-linearity and admissibility by
distributivity. Evaluation $\phi\mapsto\phi(k)$ maps it into
$K(l,k)$. Under the preceding identification, the image of a coefficient
$x\in F(\phi(k))$ is precisely the defining canonical generator
$t_{\phi(k)l,k}a_{\phi(k),x,l}$. The image of the displayed colimit
map is an additive subgroup containing all these generators and no
others beyond their sums, so it equals $\mathcal J_L(l,k)$.
Directedness combines every finite sum at one admissible functional.
No injectivity of any coefficient map is used.

If $Q$ is finite, $K(l,k)$ has the greatest element $r(l,k)$, and its
colimit is $F(r(l,k))$. The inequality
$r(h,k)r(l,h)\leq r(l,k)$ follows from its action on $l$, and yields
the stated multiplication formula.

For arbitrary $L$, the pointwise join of all admissible functionals
still exists when $Q$ is finite. At each argument there are only
finitely many distinct values. Hence finite distributivity proves that
this join preserves zero, binary joins and the left $Q$-action.
For its admissibility at a fixed $k$, select finitely many functionals
realizing those distinct values; their joined value acting on $l$ is
at most $k$. The resulting map is therefore the greatest member
$\delta_l$ of $D_l$. No finite global subfamily of functionals is
asserted. The second colimit is consequently $F(\delta_l(k))$, and
its map to the first is the coefficient transition to $F(r(l,k))$.
This proves the finite-$Q$ formulas with arbitrary $L$ and $F$.
\end{proof}
The maximum $r(l,k)$ is the scalar-action residual. For the classical
residuation convention, compare Cohen, Gaubert and Quadrat
\cite[Sections~2.1--2.2, equations~(2)--(3) and~(13)--(14) of the
preprint]{CGQ}; their displayed semimodule convention is right-sided.
Here it is the left-action residual. Completeness is not assumed in the
general theorem: the colimits replace unavailable maxima.

\begin{corollary}[ordinary group-valued modules]
\label{cor:group-projectives-general}
Let
\[
 R(T)=\varinjlim_{q\in Q}F(q),\qquad
 I_0=\im\bigl(F(0_Q)\longrightarrow R(T)\bigr).
\]
Coefficient multiplication makes $R(T)$ the universal ordinary ring
reflection of $T$. Semimodules with singleton additive support are
exactly ordinary $R(T)$-modules. Such a module $P$ is finitely generated
projective over $T$ if and only if it is finitely generated projective
over $R(T)$ and $I_0P=P$.
\end{corollary}
\begin{proof}
The poset $Q$ is directed by finite joins. Naturality and associativity
of $\mu$ make its group colimit a unital ring. The map sending
$(q,x)\in T$ to the class of $x$ preserves both operations, zero and
unit. Every semiring map from $T$ to a ring sends additive idempotents
to zero, so its restrictions to the groups $F(q)$ are compatible.
They induce a unique unital ring map from this colimit, proving the
universal property.

For the singleton-zero support $L=\{o\}$, one has $K(o,o)=Q$, so
$\mathcal A_L=R(T)$. Its only $Q$-linear functional to $Q$ is zero;
therefore $\mathcal J_L=I_0$. The fixed-support equivalence and
Theorem~\ref{thm:canonical-ideal} prove the remaining statements,
since $L$ is a finite projective $Q$-support. This includes $R(T)=0$.
\end{proof}
For $T=R_0$ the ideal $I_0$ vanishes, recovering
the one-idempotent obstruction in Section~\ref{sec:projectives}; for the two-level ideal
construction it recovers Proposition~\ref{prop:canonical-arbitrary-ideal}.
